\section{Conclusion}
\label{sec:conclusion}

This paper examined a narrow but consequential failure mode in AI-assisted cyber-physical security: an explanation can preserve the detector's committed facts while changing the operator's decision frame. In the frozen held-out protocol, a conventional structural verifier accepted 55.7\% of eligible external-model candidates and at least one of four candidates on 78.8\% of tasks. \system{} moved integrity-critical fields out of the generative trust boundary, committed them to a deterministic rendering of a closed evidence bundle, and verified that commitment independently. The hardened gate had zero observed candidate and task ASR, zero clean rejections, and sub-millisecond verification overhead in the staged CPU measurements. A separate four-model development analysis found zero registered-surface bypass among 288 eligible candidates while preserving every tested clean and preregistered benign surface.

The completed blinded review shows that 12/80 sampled attack texts meet a frozen semantic-validity rule, with high inter-reviewer reliability and independent resolution of 14 disagreements. It also exposes an important boundary: 4/24 diagnostic controls meet the same rule, so the rubric is evidence about text properties rather than a deployable classifier. The detector remains imperfect, its confidence requires better calibration, the process graph contains relations with different evidential strength, and no operator study has shown how the text changes behavior. The next phase should therefore combine newly sampled multi-model and cross-dataset red teaming with authenticated telemetry provenance, nonvacuous process invariants, and controlled operator experiments. Within the tested boundary, the evidence supports a practical design principle: generative AI may help communicate an alarm, but it should not have unilateral authority to redefine the evidence, urgency, or action shown to a human operator.
